% Preface, translated from sv/front/preface.tex. No class, year or date: the book is meant to
% outlast them.

\chapter{Preface}
\markboth{Preface}{Preface}

This book contains the mathematics used in electrical engineering. It is set from a course in
electrical engineering mathematics, lesson by lesson, and covers:
\begin{itemize}
  \item Arithmetic, algebra, and linear equations and systems of equations.
  \item Powers, roots, quadratic equations, and exponential functions and logarithms, including
    decibels.
  \item Vectors, functions and trigonometric functions.
  \item Derivatives and integrals, with applications in circuit theory.
  \item Complex numbers in rectangular, polar and Euler form, and phasors for AC calculations.
\end{itemize}

\textbf{The mathematics is applied to electrical problems throughout.} The first chapter already
uses fractions on resistors in parallel, and the book keeps returning to circuits:
\begin{itemize}
  \item Resistors in series and in parallel, voltage dividers and power calculations.
  \item Charging and discharging capacitors in RC circuits.
  \item AC voltages: amplitude, frequency, phase and trigonometric equations.
  \item Impedance calculations with complex numbers, and phasor addition.
  \item Sampling sinusoidal signals, and computing derivatives, integrals and complex numbers in
    C, in the lab in \kapref{18}.
\end{itemize}

\section*{What you will be able to do}
Once you have worked through the book, you should be able to:
\begin{itemize}
  \item Simplify algebraic expressions and solve the equations and systems of equations that
    arise in electrical engineering.
  \item Calculate with powers, roots, logarithms and decibels.
  \item Calculate and interpret derivatives and integrals of common functions.
  \item Calculate with complex numbers and phasors to analyse AC circuits.
  \item Implement the mathematics of the course in C with \file{math.h} and \file{complex.h}.
\end{itemize}

\section*{What you are expected to know}
The book assumes compulsory-school mathematics, and even most of that is revised in the first
chapters: number sets, the order of operations, fractions and percentages in \kapref{1},
equations in \kapref{3} and powers in \kapref{5}. If you feel unsure of the basics, simply start
from the beginning.

The lab in \kapref{18} also assumes basic C programming: variables, functions, \code{printf()}
and loops. Floating-point numbers, \file{math.h} and \file{complex.h} are covered in the chapter.

\section*{How the book is organised}
Each chapter corresponds to a lesson in the course and has the same number: Chapter~5 is
lesson~5, and its material is in the course repository under \file{lectures/L05}.

\begin{booktable}{@{}p{5em}L@{}}
  \thead{Chapter} & \thead{Topic} \\
  \midrule
  1 & Number systems and basic arithmetic \\
  2 & Algebra \\
  3 & Linear equations and inequalities \\
  4 & Systems of equations \\
  5 & Powers and roots \\
  6 & Quadratic equations, with practice test 1 \\
  7 & Vectors \\
  8 & Functions \\
  9 & Trigonometric functions \\
  10 & Exponential functions, logarithms and decibels \\
  11 & Review of Chapters 7--10, with practice test 2 \\
  12 & Derivatives, part I \\
  13 & Derivatives, part II \\
  14 & Integrals \\
  15 & Complex numbers, part I \\
  16 & Complex numbers, part II \\
  17 & Complex numbers, part III \\
  18 & Mathematics in C: lab \\
  19 & Review and exam preparation, with practice exam \\
\end{booktable}

A chapter opens with what it contains and then goes through the theory, with examples that show
every step. Then comes a summary with the chapter's formulas, what you should be able to do and a
few questions under the heading \emph{Test yourself}, and last the chapter's exercises.
Chapters~11 and~19 have no new theory; they review, and hold a practice test and a practice exam
respectively.

\section*{The exercises}
The exercises in each chapter are divided into parts and numbered as in the course, so that
Exercise~2.3 is the third exercise in Part~2:

\begin{booktable}{@{}p{5em}L@{}}
  \thead{Part} & \thead{What it contains} \\
  \midrule
  Part 1 & Introductory exercises, or review of the previous chapter. \\
  Part 2 & New material: the chapter's own theory, often applied to a circuit. \\
  Part 3 & Extra exercises for further practice after the lesson. \\
\end{booktable}

In the course, Parts~1 and~2 are done during the lesson, first on your own and then together,
while Part~3 is optional practice.

\section*{Answers and worked solutions}
\textbf{The answers to all the exercises are in the answer key at the back of the book.} Every
exercise has a small tag, \svarexempel, giving the page its answer is on; in the PDF, both the tag
and the exercise's number in the answer key are links between the exercise and its answer.

\textbf{The answer key holds only the answers, not the solutions.} An answer shows whether you
got it right, but not how to get there. Full worked solutions, in Swedish, are published in the
course repository after each lesson, under \file{lectures/LNN/appendix/c_solutions.md}. Always
work an exercise yourself first, then compare with the answer key, and read the worked solution
when your answer does not match or when you want to see another way.

\section*{What you need}
A calculator, any calculator, is used throughout the book. Many exercises are still meant to be
done without one, and they say so. The lab in \kapref{18} needs nothing but a web browser: the
code is written and run in an online compiler, with nothing to install.

\section*{Notation}
The book follows the notation common in electrical engineering, and that of the Swedish edition it
is translated from:
\begin{itemize}
  \item \textbf{Decimal point}: $4.7$. The Swedish edition, the course material and the exams
    write a decimal comma, $4{,}7$, for the same number.
  \item \textbf{Units} are written with a small space after the number: $4.7\,\text{k}\Omega$,
    $50\,\text{Hz}$, $30\,\%$.
  \item \textbf{Voltage} is written $U$, as in international standards and in the Swedish course;
    many English-language texts write $V$ instead. Subscripts are English: $U_{\text{out}}$ is
    the output voltage.
  \item \textbf{Resistance is written before current}, as in Ohm's law: $U = RI$ and $P = RI^2$.
  \item \textbf{Parallel connection} is written $R_1 \parallell R_2$.
  \item \textbf{Logarithms}: $\log$ is the base-10 logarithm and $\ln$ the natural logarithm.
  \item \textbf{The imaginary unit} is written $j$, since $i$ already stands for current.
  \item \textbf{Angles} are in radians unless otherwise stated; degrees are marked with $^{\circ}$.
\end{itemize}
